\begin{abstract}
We attribute $512\times512$ JPEGs to a camera or to one of two undisclosed generators with a single from-scratch model of at most 12.3M parameters. A ResNet-18 that reads the image and its neighbouring-pixel residual at full resolution scores 100\% on held-out and development data, yet labels all 1,500 test images as real. From the unlabelled test spectra we infer why: the test images were stretched by about $512/508$ and re-encoded as JPEG q90 4:2:0. The stretch shifts the generators' period-8 up-sampling lattice by up to 4 pixels, and the detector is locked to that lattice's phase. Our submitted model, v4, keeps the residual input but down-samples with anti-aliasing BlurPool layers. It is trained on $256^2$ crops of 152k stored, broadly augmented copies and applied once to each test image as given. It scores \textbf{95.3\%} on test (the baseline's all-real output scores 33.3\%) and 96.1\% on development data. Probes show that v4 is invariant to sub-period shifts and flips. An ablation shows that augmentation narrowly matched to the inferred processing overfits it. Almost all remaining errors are real photographs labelled as fake.
\end{abstract}
